\documentclass[12pt]{extarticle}

% ===== Encodage & langue =====
\usepackage{fontspec}
\usepackage{polyglossia}
\setdefaultlanguage{french}

% ===== Police Verdana (compiler avec XeLaTeX ou LuaLaTeX) =====
\setsansfont{Verdana}
\renewcommand{\familydefault}{\sfdefault}

% ===== Interligne =====
\usepackage{setspace}
\setstretch{1.1}

% ===== Marges réduites =====
\usepackage{geometry}
\geometry{a4paper, top=1.8cm, bottom=1.8cm, left=2cm, right=2cm}

% ===== Mathématiques =====
\usepackage{amsmath}

% ===== Espacement réduit autour des équations =====
\setlength{\abovedisplayskip}{4pt}
\setlength{\belowdisplayskip}{4pt}
\setlength{\abovedisplayshortskip}{2pt}
\setlength{\belowdisplayshortskip}{2pt}

% ===== Espacement réduit des sections =====
\usepackage{titlesec}
\titlespacing*{\section}{0pt}{1.2ex}{0.6ex}

% ===== Titre du document =====
\title{\textbf{Correction du contrôle n\textsuperscript{o} 1}}
\author{}
\date{}

\begin{document}
\maketitle

% --------------------------------------------------
\section*{Exercice 1}
\begin{align*}
18 - 6 - 2 + 19 - 3 &= 12 - 2 + 19 - 3 \\
                     &= 10 + 19 - 3 \\
                     &= 29 - 3 \\
                     &= 26
\end{align*}

% --------------------------------------------------
\section*{Exercice 2}

\begin{align*}
\textbf{a)}
5 + 3 \times 2 &= 5 + 6 \\
               &= 11
\end{align*}

\textbf{b)}
\begin{align*}
34 - 6 \times 2 &= 34 - 12 \\
                &= 22
\end{align*}

\textbf{c)}
\begin{align*}
29 + 12 \div 4 &= 29 + 3 \\
               &= 32
\end{align*}

\textbf{d)}
\begin{align*}
6 + 9 \times 3 + 2 - 4 \div 4 &= 6 + 27 + 2 - 1 \\
                              &= 34
\end{align*}

% --------------------------------------------------
\section*{Exercice 3}
\begin{align*}
3 \times 7 + 2 \times 8 &= 21 + 16 \\
                        &= 37
\end{align*}

La dépense est de 37~€.

% --------------------------------------------------
\section*{Exercice 4}

\textbf{a)} $27 - (6 + 4) \div 2$
\begin{align*}
&= 27 - 10 \div 2 \\
&= 27 - 5 \\
&= 22
\end{align*}

\textbf{b)} $(3 + 8) \times 3 - 7 + 5$
\begin{align*}
&= 11 \times 3 - 7 + 5 \\
&= 33 - 7 + 5 \\
&= 31
\end{align*}

\textbf{c)} $(19 - (5 + 7)) \times 3$
\begin{align*}
&= (19 - 12) \times 3 \\
&= 7 \times 3 \\
&= 21
\end{align*}

\textbf{d)} $4 \times ((21 - (6 + 13)) - 7$
\begin{align*}
&= 4 \times (21 - 19) - 7 \\
&= 4 \times 2 - 7 \\
&= 8 - 7 \\
&= 1
\end{align*}

% --------------------------------------------------
\section*{Exercice 5}
\begin{align*}
11 - 4 \times 2 &= 11 - 8 \\
                &= 3
\end{align*}

Une gomme coûte 3~€.

% --------------------------------------------------
\section*{Exercice bonus 1}

\textbf{a)} $7 \times 3 + 13$

\textbf{b)} $9 \times (7 + 1)$

% --------------------------------------------------
\section*{Exercice bonus 2}

\textbf{a)} La somme du produit de 8 par 3 et de 7.

\textbf{b)} Le produit de la somme de 6 et de 5 par 4.

\end{document}